\begin{tikzpicture}[font=\small]
  % ---------- right rotation ----------
  \begin{scope}
    \node[anchor=south, font=\small] at (2.6,3.2){右旋（Right Rotation）：左子树过高时用};

    % before
    \node[bad]  (y) at (0.9,2.3) {$y$};
    \node[vtx] (x) at (0.0,1.3) {$x$};
    \node[vtx, minimum size=6mm, font=\footnotesize] (C) at (1.8,1.3) {$C$};
    \node[vtx, minimum size=6mm, font=\footnotesize] (A) at (-0.7,0.3) {$A$};
    \node[vtx, minimum size=6mm, font=\footnotesize] (B) at (0.7,0.3) {$B$};
    \draw[edge] (y) -- (x); \draw[edge] (y) -- (C);
    \draw[edge] (x) -- (A); \draw[edge] (x) -- (B);
    \node[note, anchor=north] at (0.6,-0.25){左重：$h(x)-h(C)=2$};

    \node[note] at (3.2,1.3){$\Longrightarrow$};

    % after
    \node[vtx] (x2) at (5.0,2.3) {$x$};
    \node[vtx, minimum size=6mm, font=\footnotesize] (A2) at (4.1,1.3) {$A$};
    \node[vtx] (y2) at (5.9,1.3) {$y$};
    \node[vtx, minimum size=6mm, font=\footnotesize] (B2) at (5.2,0.3) {$B$};
    \node[vtx, minimum size=6mm, font=\footnotesize] (C2) at (6.6,0.3) {$C$};
    \draw[edge] (x2) -- (A2); \draw[edge] (x2) -- (y2);
    \draw[edge] (y2) -- (B2); \draw[edge] (y2) -- (C2);
    \node[note, anchor=north] at (5.4,-0.25){高度降 1，中序不变};
  \end{scope}

  % ---------- the four cases ----------
  \begin{scope}[yshift=-4.0cm]
    \node[anchor=south, font=\small] at (6.4,2.2){四种失衡与对应的修复};

    \foreach \i/\case/\fix in {
      0/{LL\\ 左孩子的左边过高}/{右旋一次},
      3.3/{RR\\ 右孩子的右边过高}/{左旋一次},
      6.6/{LR\\ 左孩子的右边过高}/{先对左孩子左旋\\ 再整体右旋},
      9.9/{RL\\ 右孩子的左边过高}/{先对右孩子右旋\\ 再整体左旋}}{
      \node[blk, minimum width=2.7cm, minimum height=1.0cm, anchor=north west,
            draw=algblue, fill=algblue!7] at (\i,1.7) {\case};
      \node[note, anchor=north, align=center] at (\i+1.35,0.5){\fix};
    }

    \node[note, anchor=north, align=center] at (6.4,-1.0)
      {记忆线索：\textbf{同向}（LL / RR）转一次，\textbf{异向}（LR / RL）先掰直再转。\\
       每次插入/删除后沿着回溯路径检查，最多 $O(\log n)$ 次旋转；插入其实最多 \textbf{1} 次。};
  \end{scope}

  % ---------- why height stays O(log n) ----------
  \node[note, anchor=north west, align=left] at (0,-7.6)
    {\textbf{为什么高度是 $O(\log n)$}：设高度为 $h$ 的 AVL 树最少有 $N(h)$ 个结点，则\\
     $N(h)=1+N(h-1)+N(h-2)$ —— 这正是斐波那契的形状，$N(h)$ 随 $h$ \textbf{指数增长}。\\
     反过来 $h=O(\log n)$。「每个结点左右高度差不超过 1」这条局部约束，逼出了全局的对数高度。};
\end{tikzpicture}
